\section*{集合论。}
\phantomsection
\addcontentsline{toc}{section}{集合论。}

可以说，在任何时代，数学家和哲学家都在或多或少自觉地使用着来自集合论的推理；但在他们关于这一主题的概念的历史中，必须把一切与基数观念（特别地，与无穷观念）相联系的问题，同那些仅仅引入属于与包含这两个概念的问题截然分开。后者属于最直观的概念之列，似乎从未引起过争议：正是可以在它们之上最方便地建立三段论的理论（如莱布尼茨和欧拉将向我们表明的那样），或者诸如“整体大于其各部分之和”之类的公理，更不必说几何中处理曲线与曲面相交的内容了。直到 19 世纪末，谈论由具有这样或那样给定性质的对象组成的集合（某些作者称之为“类”）都不曾出现任何问题；\(^{38}\) 而康托尔给出的那个著名的“定义”（“所谓集合，是指把我们的直觉或思维中彼此分明的一些对象汇集为一个整体”（{[}47{]}，第~282~页））在其发表之时几乎没有引起任何异议。\(^{39}\) 然而一旦集合的概念与数或量的概念掺和在一起，情况就完全不同了。直线的无限可分性问题（最早的一批毕达哥拉斯学派学者无疑已经提出了它）众所周知后来引出了相当大的哲学困难：从埃利亚学派到波尔查诺和康托尔，数学家和哲学家们前仆后继地扑向“由无穷多个无大小的点组成的有限的量”这一悖论，却均未成功。即使概略地重述这一问题所激起的无休止而激烈的论战——它们构成一块特别有利于形而上学或神学离题发挥的土壤——对我们来说也没有任何益处；我们只需指出自古以来大多数数学家所止步的那个观点。这一观点实质上在于：只要还不能以无可辩驳的方式作出结论，就拒绝讨论——我们将在现代形式主义者那里重新见到这种态度；正如后者设法排除“悖论的”集合的介入一样（见下文第~31~页以下），经典数学家也小心翼翼地避免在推理中引入“实无穷”（也就是说，避免引入由被设想为同时存在的、至少在思想中同时存在的无穷多个对象所组成的集合），而满足于“潜无穷”，即增大每个给定量的可能性（若是“连续”的量，则是减小它的可能性）。\(^{40}\) 尽管这一观点含有一定分量的虚伪，\(^{41}\) 它却毕竟容许了经典数学的绝大部分（包括比例理论以及后来的无穷小演算）得到发展；\(^{42}\) 尤其是在无穷小所引起的论战之后，它甚至被视为一种绝妙的保险，直到深入 19 世纪之后很久，它几乎仍是人们普遍信奉的一条教条。

关于等势这一一般概念的第一缕微光，出现在伽利略的一段评注中（{[}122 b{]}，第 VIII 卷，第~78-80~页）：他注意到映射 \(n \mapsto n^2\) 在自然数与其平方数之间建立了一个双射对应，并由此得出结论：公理“整体大于部分”

不能应用于无穷集合。但是，这一评注远没有开创对无穷集合的理性研究，它似乎除了加深人们对实无穷的不信任之外别无其他效果；这已经是伽利略本人所得出的结论，而柯西在 1833 年引用他，也只是为了赞同他的这一态度。

正是分析的需要——特别是贯穿整个 19 世纪的对实变量函数的深入研究——构成了后来成为现代集合论的那门理论的源头。当波尔查诺在 1817 年证明 R 中一个有下界的集合的下极限存在时（{[}27 c{]}），他仍像他同时代的大多数人那样“按内涵”推理，谈论的不是任意的实数集合，而是实数的任意性质。但 30 年后，当他撰写《无穷的悖论》（Paradoxien des Unendlichen）{[}27 b{]}（1851 年出版，即他去世三年之后）时，他已毫不犹豫地主张“实无穷”存在的权利，并谈论任意的集合了。他在这部著作中定义了两个集合等势的一般概念，并证明了 R 中两个紧区间是等势的；他还注意到有限集与无穷集的特征区别在于：无穷集 E 与一个异于 E 的子集等势；但他对这一断言没有给出任何令人信服的证明。这部著作的总基调终究是哲学多于数学，而且由于未能把集合的势的概念同无穷的大小或等级的概念足够清楚地分离开来，波尔查诺构造势越来越大的无穷集的尝试均告失败，他还借此机会把推理与关于发散级数的、完全没有任何意义的考虑混在了一起。

我们今天所知的集合论之得以创立，应归功于 G. 康托尔的天才。他同样是从分析出发的：他关于三角级数的工作（受黎曼工作的启发）很自然地使他在 1872 年对出现在这一理论中的“例外”集合作出第一次分类尝试，\(^{43}\) 所借助的则是他在此时引入的逐次“导集”的概念。无疑正是与这些研究、也与他定义实数的方法（{[}47{]}，第~92-97~页）相联系，康托尔开始对等势问题发生兴趣，因为在 1873 年他指出，有理数集（或代数数集）是可数的；而在大约始于此时与戴德金的通信 {[}48{]} 中，我们看到他提出了整数集与实数集之间的等势问题，几周之后他便成功地给出了否定的解决。随后，从 1874 年起，占据他心灵的是维数问题：有三年的时间，他徒劳地试图证明 R 与 \(R^{n}(n > 1)\) 之间不可能存在双射对应，最终却大出其所料地\(^{44}\)定义出了这样一个对应。有了这些既新奇又惊人的结果，他便把自己完全献身于集合论。在 1878 年至 1884 年间发表于《数学年刊》（Mathematische Annalen）的一系列 6 篇论文中，他同时处理了等势问题、全序集理论、R 与 R \(^{n}\) 的拓扑性质以及测度问题；令人赞叹的是，可以看到这些在经典“连续统”观念看来如此错综纠缠的概念，在他手中是如何一点一点地清晰呈现出来的。从 1880 年起，他有了“超限地”迭代“导集”构造的想法；但这一想法直到两年之后，随着良序集——康托尔最独创的发现之一——的引入才获得生命；多亏良序集，他才能够着手对基数作详细的研究，并提出“连续统问题”{[}47{]}。

这样大胆的观念推翻了两千年的传统，又导出如此出人意料、外表如此悖理的结果，它们不可能不经抵抗就被接受。事实上，在当时德国有影响的数学家中，只有魏尔斯特拉斯一人以几分好感关注着康托尔（他昔日学生）的工作；另一方面，康托尔却遇到了施瓦茨尤其是克罗内克的不可调和的反对。\(^{45}\) 看来，使康托尔出现神经疾病的最初症状、并使他的数学创造力因之受损的，既是由对其思想的反对所造成的持续紧张，同样也是他为证明连续统假设而付出的徒劳努力。\(^{46}\) 他直到 1887 年才真正重新拾起对集合论的兴趣，他最后的论著发表于 1895-1897 年；其中他主要发展了全序集理论和序数的演算。他还在 1890 年证明了不等式 \(m < 2^{m}\)；然而，不仅连续统问题仍无答案，而且在基数论中还留下一个更严重的空缺：康托尔未能确立任意基数之间良序的存在。这一空缺即将得到填补：一方面靠 F. 伯恩斯坦定理（1897），它表明关系 \(a \leq b\) 与 \(b \leq a\) 蕴涵 \(a = b\)，\(^{47}\) 另一方面尤其是靠策梅洛定理 {[}342 a{]}，它证明了每个集合上都存在良序——这个定理康托尔早在 1883 年就已猜测到了（{[}47{]}，第~169~页）。

然而戴德金从一开始便以持续的关注追随着康托尔的研究；不过，当后者把注意力集中于无穷集合及其分类时，戴德金却循着他自己关于数概念的思路（这条思路此前已使他借“分割”给出了无理数的定义）。在他 1888 年出版的小册子《数是什么，应是什么》（Was sind und was sollen die Zahlen）中——其要旨早在 1872-78 年即已写成（{[}79{]}，第 III 卷，第~335~页）——他表明了自然数概念（如前所见，整个经典数学最终都依赖于它）本身如何能够从集合论的基本概念中导出。他（无疑是最早明确地这样做的人）发展了一个集合到另一个集合的任意映射的初等性质（这在康托尔之前一直无人问津，康托尔只对双射对应感兴趣）；对于集合 E 到自身的每个映射 f，他引入了元素 \(a \in E\) 相对于 f 的“链”的概念，即集合 \(K \subset E\) 中所有满足 \(a \in K\) 且 \(f(K) \subset K\) 者的交。⁴⁸ 随后，他把无穷集 E 定义为：存在 E 到 E 内的一个双射 \(\phi\) 使得 \(\phi(E) \neq E\)；⁴⁹ 进一步，若还存在这样一个映射 \(\phi\) 和一个元素 \(a \notin \phi(E)\)，使得 E 是 a 的链，戴德金便称 E 是“单纯无穷的”；他注意到此时“皮亚诺公理”即得到满足，并（先于皮亚诺）表明了如何从出发得到初等算术的全部定理。他的阐述中唯一缺失的部分是无穷公理：戴德金（步波尔查诺的后尘）自认为通过把人的“思想世界”（“Gedankenwelt”）看作一个集合就能证明这条公理。⁵⁰

从另一条路走，戴德金由于其在算术方面的工作（特别是理想理论）而被引去在比康托尔更一般的框架中考虑有序集的概念。当后者把自己完全局限于全序集时，\(^{51}\) 戴德金则处理一般的情形，并对格作了特别深入的研究（{[}79{]}，第 II 卷，第~236-271~页）。这项工作在当时几乎无人注意；尽管这些被好几位作者重新发现的结果自 1935 年以来成为大量出版物的对象，它们的历史重要性与其说来自这一理论无疑相当有限的应用可能，不如说来自如下事实：它们构成了精心进行公理化构造的最早例子之一。另一方面，康托尔关于可数集或具有连续统的势的集合的最初结果，很快就有了为数众多而且重要的应用，直至延伸到分析中最经典的问题\(^{52}\)（自然还不必说康托尔工作中开创了一般拓扑学和测度论的那些部分；关于这一点见第~141~页和第~221~页）。此外，从 19 世纪的最后几年起，超限归纳原理的最初应用开始出现；尤其在策梅洛定理获证之后，它已成为现代数学所有分支中不可或缺的工具。库拉托夫斯基于 1922 年给出了这条原理的一个往往更合用的版本，它避免了良序集的使用（{[}189 a{]}，第~89~页）；今天人们主要使用的正是这个形式，它后来又由佐恩重新发现 {[}344{]}。\(^{53}\)

到 19 世纪末，康托尔的基本概念就这样站住了脚。\(^{54}\) 我们已经看到，大约在这个时候，数学的形式化宣告完成，公理化方法的使用也几乎取得了一致的同意。但与此同时，一场罕见的剧烈的“基础危机”开启了，它震撼数学界达三十余年，而且有时看来不仅危及一切新近的收获，甚至危及数学中最经典的部分。

